\section{SECCIÓN 5. ESTÁNDARES, PRÁCTICAS, CONVENCIONES Y MÉTRICAS}
Para verificar la entrega de un producto totalmente conforme y de alta calidad, cada individuo asignado al proyecto participará en el aseguramiento de la calidad. La referencia (e) define los procedimientos mediante los cuales el personal de desarrollo de software verificará la calidad del producto durante el proceso de desarrollo. El resto de esta sección describe los procedimientos utilizados por SQA para verificar que se cumplan las disposiciones de aseguramiento de la calidad de este Plan SQA y los estándares, prácticas, convenciones y métricas aplicables.

El presente Plan SQA define los estándares, prácticas, convenciones y métricas aplicables al proyecto Galería de Imágenes Segura. La Sección 3 identifica la evaluación de SQA de las fases de implementación y prueba para verificar el cumplimiento de estos estándares.

La Sección 4 identifica los Configuration Items (CIs) asociados a cada producto de software que se desarrollará. SQA verificará que el formato y el contenido de la documentación cumplan con los criterios definidos en el SQAP.

Los estándares para la estructura lógica, la codificación y los comentarios del código se describen en las convenciones del proyecto (Google Java Style Guide para backend, Airbnb Style Guide para frontend). SQA verificará que el código fuente cumpla con estos estándares.

Los estándares y prácticas para las pruebas se describen en el Plan de Pruebas de Software (STP, CI-04). SQA verificará que las actividades de prueba cumplan con lo establecido.

\subsection{5.1 MÉTRICAS}
\textit{Identificar o hacer referencia a los estándares, prácticas y convenciones que se utilizarán en la definición, recopilación y utilización de datos de medición de software. Citar cualquier requisito o estándar interno (ej., del proyecto, corporativo) y externo (ej., del usuario, cliente) con el que deban cumplir las prácticas de métricas. IEEE Std 1045-1992, IEEE Standard for Software Productivity Metrics, referencia (p), describe convenciones para contar los resultados de los procesos de desarrollo. IEEE Std 1061-1992, IEEE Standard for a Software Quality Metrics Methodology, referencia (q), proporciona una metodología para seleccionar e implementar métricas de proceso y producto. IEEE Std 982.1-1988, IEEE Standard Dictionary of Measures to Produce Reliable Software, referencia (r), e IEEE Std 982.2-1988, IEEE Guide for using reference (r), referencia (s), proporcionan varias medidas para su uso en diferentes fases del ciclo de vida para ganar confianza en la construcción de software confiable. Para mantener las métricas simples, se ofrece un ejemplo de métricas de costo y cronograma.}

Las siguientes mediciones se realizarán y utilizarán para determinar el estado de costo y cronograma de las actividades de SQA:
\begin{enumerate}
    \item Fechas de hitos de SQA (planificadas)
    \item Fechas de hitos de SQA (completadas)
    \item Trabajo de SQA programado (planificado)
    \item Trabajo de SQA completado (real)
    \item Esfuerzo de SQA dedicado (planificado)
    \item Esfuerzo de SQA dedicado (real)
    \item Fondos de SQA gastados (planificados)
    \item Fondos de SQA gastados (reales)
    \item Número de elementos de incumplimiento abiertos
    \item Número de elementos de incumplimiento cerrados
    \item Número total de elementos de incumplimiento
\end{enumerate}

SQA es responsable de reportar estas mediciones al Gerente de Proyecto mensualmente.

\subsection{5.2 MÉTRICAS DE CALIDAD EN USO }

La norma ISO/IEC 25022 define métricas para medir la calidad en uso de un producto de software desde la perspectiva del usuario final. Las métricas presentadas a continuación se organizan según las cinco características definidas en dicha norma: efectividad, eficiencia, satisfacción, libertad de riesgo y cobertura de contexto.

\subsubsection{Efectividad}

Precisión e integridad con las que los usuarios alcanzan objetivos específicos.

\begin{enumerate}
    \item \textbf{Tasa de finalización de tareas (Task completion rate)}
          \[
              TCR = \frac{N_{tareas\_completadas}}{N_{tareas\_intentadas}} \times 100\%
          \]
          Donde \(N_{tareas\_completadas}\) es el número de tareas (pruebas, auditorías) completadas exitosamente y \(N_{tareas\_intentadas}\) es el total de tareas ejecutadas.

    \item \textbf{Frecuencia de errores (Error frequency)}
          \[
              FE = \frac{N_{errores}}{N_{ejecuciones}}
          \]
          Donde \(N_{errores}\) es el número de ejecuciones que presentaron errores y \(N_{ejecuciones}\) es el número total de ejecuciones realizadas.

    \item \textbf{Tasa de logros de efectividad (Effectiveness achievement rate)}
          \[
              EAR = \frac{N_{audits\_aprobados}}{N_{audits\_totales}}
          \]
          Donde \(N_{audits\_aprobados}\) son las auditorías que cumplen el umbral definido y \(N_{audits\_totales}\) es el total de auditorías ejecutadas.

    \item \textbf{Calidad percibida del producto}
          \[
              CP = \frac{\sum_{i=1}^{n} w_i \cdot s_i}{\sum_{i=1}^{n} w_i}
          \]
          Donde \(w_i\) es el peso asignado a cada categoría de calidad y \(s_i\) es el puntaje obtenido en esa categoría (escala 0--1).
\end{enumerate}

\subsubsection{Eficiencia}

Recursos gastados en relación con la precisión e integridad con las que los usuarios alcanzan los objetivos.

\begin{enumerate}
    \item \textbf{Tiempo de tarea (Task time)}
          \[
              TT = t_{fin} - t_{inicio}
          \]
          Medición directa del tiempo transcurrido desde el inicio hasta la finalización de una tarea, expresado en milisegundos.

    \item \textbf{Índice de eficiencia de carga (Load efficiency index)}
          \[
              LEI = \frac{LCP}{TTI}
          \]
          Donde \(LCP\) es el Largest Contentful Paint y \(TTI\) es el Time to Interactive. Un valor cercano a 1 indica carga continua sin bloqueos.

    \item \textbf{Tiempo de respuesta del sistema}
          \[
              TR = RTT + L_{servidor}
          \]
          Donde \(RTT\) es el Round Trip Time de red y \(L_{servidor}\) es la latencia del servidor backend.

    \item \textbf{Peso de recurso por elemento DOM}
          \[
              PR = \frac{B_{totales}}{N_{DOM}}
          \]
          Donde \(B_{totales}\) es el peso total en bytes de los recursos cargados y \(N_{DOM}\) es el número de elementos en el DOM.

    \item \textbf{Proporción de tiempo de ejecución JavaScript}
          \[
              PJS = \frac{T_{bootup}}{T_{mainthread}} \times 100\%
          \]
          Donde \(T_{bootup}\) es el tiempo de ejecución de JavaScript y \(T_{mainthread}\) es el tiempo total de trabajo del hilo principal.
\end{enumerate}

\subsubsection{Satisfacción}

Grado en que las necesidades del usuario se satisfacen cuando se utiliza el producto.

\begin{enumerate}
    \item \textbf{Puntaje de satisfacción del usuario (User satisfaction score)}
          \[
              USS = \frac{\sum_{i=1}^{n} R_i}{n}
          \]
          Donde \(R_i\) son las calificaciones individuales obtenidas mediante encuestas (escala Likert 1--5) y \(n\) es el número de encuestados.

    \item \textbf{Net Promoter Score (NPS)}
          \[
              NPS = P_{promotores} - P_{detractores}
          \]
          Donde \(P_{promotores}\) es el porcentaje de usuarios que califican \(9\)--\(10\) y \(P_{detractores}\) el porcentaje que califica \(0\)--\(6\).

    \item \textbf{Tasa de utilidad percibida (Perceived usefulness)}
          \[
              UP = \frac{N_{funcionalidades\_utilizadas}}{N_{funcionalidades\_disponibles}} \times 100\%
          \]
\end{enumerate}

\textit{Nota: Las métricas de satisfacción requieren la aplicación de instrumentos de recolección como encuestas SUS (System Usability Scale) o cuestionarios personalizados. Se planifica su implementación en fases posteriores.}

\subsubsection{Libertad de Riesgo}

Grado en que un producto mitiga los riesgos económicos, de salud, seguridad o medio ambiente.

\begin{enumerate}
    \item \textbf{Tasa de mitigación de riesgos de seguridad (Security risk mitigation)}
          \[
              SRM = 1 - \frac{V_{criticas} \cdot 10 + V_{altas} \cdot 5 + V_{medias} \cdot 2}{N_{endpoints}}
          \]
          Donde \(V_{criticas}\), \(V_{altas}\) y \(V_{medias}\) son vulnerabilidades encontradas según su severidad y \(N_{endpoints}\) es el número de endpoints auditados.

    \item \textbf{Índice de riesgo frontend}
          \[
              IRF = 1 - BP
          \]
          Donde \(BP\) es el puntaje de Best Practices obtenido en Lighthouse (escala 0--1).

    \item \textbf{Vulnerabilidad operativa detectada}
          \[
              VO = \frac{N_{errores\_consola}}{N_{paginas\_auditadas}}
          \]
          Donde \(N_{errores\_consola}\) es el número de páginas que presentan errores en la consola del navegador.

    \item \textbf{Densidad de vulnerabilidades en dependencias}
          \[
              DVD = \frac{N_{vulnerabilidades\_npm}}{N_{dependencias}}
          \]
          Donde \(N_{vulnerabilidades\_npm}\) es el conteo de vulnerabilidades reportadas por \texttt{npm audit} y \(N_{dependencias}\) es el total de dependencias del proyecto frontend.
\end{enumerate}

\subsubsection{Cobertura de Contexto}

Grado en que un producto puede ser utilizado en contextos específicos o más allá de los inicialmente previstos.

\begin{enumerate}
    \item \textbf{Completitud de cobertura de pruebas}
          \[
              CCP = \frac{M_{cubiertos}}{M_{totales}} \times 100\%
          \]
          Donde \(M_{cubiertos}\) son los módulos o capas del sistema que tienen pruebas automatizadas y \(M_{totales}\) es el total de módulos identificados (unitarias, integración, E2E, carga, seguridad).

    \item \textbf{Madurez de aseguramiento de calidad}
          \[
              MA = \frac{C_{ejecutadas}}{C_{totales}} \times 100\%
          \]
          Donde \(C_{ejecutadas}\) son las capas de prueba con resultados disponibles y \(C_{totales}\) son las capas definidas en el plan de pruebas.

    \item \textbf{Índice de accesibilidad}
          \[
              IA = A
          \]
          Donde \(A\) es el puntaje de accesibilidad obtenido en Lighthouse (escala 0--1), que refleja la capacidad del producto para ser utilizado por personas con discapacidades.

    \item \textbf{Tasa de adaptabilidad a dispositivos}
          \[
              TAD = \frac{N_{viewport\_configurados}}{N_{pantallas\_soportadas}}
          \]
          Donde \(N_{viewport\_configurados}\) son los puntos de ruptura responsivos implementados y \(N_{pantallas\_soportadas}\) son las resoluciones objetivo definidas en los requisitos.
\end{enumerate}
\subsection{Métricas de Calidad de Producto}
La norma ISO/IEC 25023 define métricas para medir la calidad de un producto o sistema de software de acuerdo con las características y sub-características de la ISO 25010. Las métricas presentadas a continuación se organizan según algunas de las características escogidas para medirse en este sistema: Eficiencia de desempeño, Seguridad y Mantenibilidad.

\subsubsection{Eficiencia de rendimiento}

Medidas centradas en el rendimiento y uso de recursos del sistema.

\begin{itemize}
    \item \textbf{Tiempo medio por petición}
          \[
              T_{med} = \frac{\sum t_{peticiones}}{N_{endpoints\_probados}}
          \]
          Donde $\sum t_{peticiones}$ es la suma de los tiempos registrados para las peticiones a los endpoints y $N_{endpoints\_probados}$ es el número de endpoints evaluados.

    \item \textbf{Cantidad de recursos usados}
          \[
              R_{avg} = \frac{\sum R_{peticiones}}{N_{endpoints\_probados}}
          \]
          Donde $\sum R_{peticiones}$ es la suma de los recursos consumidos (memoria, CPU u otras métricas normalizadas) por las peticiones y $N_{endpoints\_probados}$ es el número de endpoints evaluados.
\end{itemize}

\subsubsection{Seguridad}

Métricas que cuantifican la integridad y confidencialidad del sistema.

\begin{itemize}
    \item \textbf{Nivel de integridad del sistema}
          \[
              I_{nivel} = 100 - \frac{N_{amenazas\_integridad}}{N_{amenazas\_totales}} \times 100
          \]
          Donde $N_{amenazas\_integridad}$ es el número de amenazas que afectan la integridad y $N_{amenazas\_totales}$ es el total de amenazas identificadas.

    \item \textbf{Nivel de confidencialidad de datos}
          \[
              C_{nivel} = 100 - \frac{N_{amenazas\_confidencialidad}}{N_{amenazas\_totales}} \times 100
          \]
          Donde $N_{amenazas\_confidencialidad}$ es el número de amenazas que afectan la confidencialidad.
\end{itemize}

\subsubsection{Mantenibilidad}

Indicadores relativos a la facilidad de mantenimiento y calidad del código.

\begin{itemize}
    \item \textbf{Porcentaje de código duplicado}
          \[
              CD = \frac{LOC_{duplicadas}}{LOC_{revisadas}} \times 100\%
          \]

    \item \textbf{Deuda técnica promedio por módulo}
          \[
              DT_{avg} = \frac{\sum T_{deuda\_tecnica}}{N_{modulos\_revisados}}
          \]
          Donde $\sum T_{deuda\_tecnica}$ es la suma de las estimaciones (en horas o puntos) de la deuda técnica registrada por módulo.

    \item \textbf{Complejidad del código}
          \[
              CC_{avg} = \frac{\sum CC_{ciclomatica}}{N_{modulos\_revisados}}
          \]
          Donde $\sum CC_{ciclomatica}$ es la suma de la complejidad ciclomática medida en los módulos revisados.
\end{itemize}

